%==============================================================================
% PARTIDA 02: SALIDAS PARA ALUMBRADO, TOMACORRIENTES, FUERZA
% Codigos: OE.5.2.1.1.1, OE.5.2.1.1.2, OE.5.2.1.2.1, OE.5.2.1.2.2, OE.5.2.1.2.3
%==============================================================================
\subsection{OE.5.2.1.1.1 Caja Octogonal FG 100x50 mm - Salida de Alumbrado de Techo}

\subsubsection{Especificaciones Tecnicas}

La salida de alumbrado de techo consiste en la instalacion de una caja octogonal de fierro galvanizado (FG) de 100 mm x 50 mm, destinada a alojar las conexiones de la luminaria de techo. La caja debe estar provista de orejas para la fijacion de la luminaria y tapa ciega metalica durante la construccion. Se instalara empotrada en la losa de techo o en la estructura de cielo raso, segun lo indicado en los planos IE-02, IE-03 e IE-04.

La caja octogonal debe cumplir con la NTP correspondiente para cajas de derivacion metalicas, con proteccion anticorrosiva mediante galvanizado. La profundidad de 50 mm garantiza el espacio suficiente para las conexiones de conductores y la instalacion del soporte de la luminaria.

\subsubsection{Requisitos de Materiales}

\begin{itemize}
  \item Caja octogonal FG 100 mm x 50 mm con orejas para luminaria.
  \item Tapa ciega metalica para proteccion durante construccion.
  \item Conectores tipo "romex" o prensaestopas para ingreso de tuberia PVC a la caja.
  \item Tornillos de fijacion de 1/2'' para anclaje a estructura de techo.
  \item Cinta aislante vinilica para proteccion de conexiones.
\end{itemize}

\subsubsection{Metodo de Ejecucion}

\begin{enumerate}
  \item Replantear la ubicacion de la caja octogonal en el techo segun el plano de alumbrado, verificando que coincida con la posicion de la luminaria disenada.
  \item Fijar la caja octogonal a la estructura de la losa o al cielorraso mediante tornillos de anclaje, asegurando que quede firme y nivelada con la superficie del techo terminado.
  \item Conectar las tuberias de ingreso a la caja mediante conectores tipo "romex", sellando las uniones para evitar el ingreso de concreto o mortero durante el vaciado.
  \item Dejar una longitud de conductor de 0.30 m desde el borde de la caja para facilitar las conexiones futuras de la luminaria.
  \item Colocar la tapa ciega metalica para proteger la caja y los conductores durante la etapa de construccion.
  \item Identificar la salida en el tablero de distribucion correspondiente segun el circuito asignado (C1, C3 o C5).
\end{enumerate}

\subsubsection{Control de Calidad}

\begin{itemize}
  \item Verificar que la caja octogonal este correctamente nivelada y fijada.
  \item Comprobar que los conductores esten adecuadamente protegidos dentro de la tuberia hasta la caja.
  \item Inspeccionar que no existan obstrucciones en la caja que impidan la conexion futura de la luminaria.
  \item Verificar que la tapa ciega este correctamente colocada.
  \item Confirmar la ubicacion de la salida contra el plano aprobado.
\end{itemize}

\subsubsection{Metodo de Medicion}

Se medira por pieza (pza) de caja octogonal de alumbrado de techo completamente instalada, incluyendo la fijacion, conexion de tuberias, dejado de conductores y colocacion de tapa ciega.

\subsubsection{Unidad de Medida}

pza (pieza).

\subsubsection{Forma de Pago}

El pago se efectuara por pieza instalada y verificada, aprobada por la supervision. Incluye el suministro de caja, accesorios, mano de obra y herramientas necesarias para su instalacion.

%------------------------------------------------------------------------------
\subsection{OE.5.2.1.1.2 Caja Rectangular FG 4''x2'' - Salida para Interruptores}

\subsubsection{Especificaciones Tecnicas}

La salida para interruptores consiste en la instalacion de una caja rectangular de fierro galvanizado (FG) de 4 pulgadas x 2 pulgadas, destinada a alojar interruptores de luz de tipo unipolar o de conmutacion (3 vias). La caja se instalara empotrada en muros de albañileria o drywall, a una altura de 1.20 m sobre el nivel del piso terminado, medida desde el borde superior de la caja, segun el plano IE-01.

La caja debe contar con orejas de fijacion para el mecanismo del interruptor y permitir el ingreso de tuberia PVC de 3/4'' desde la caja de derivacion mas cercana.

\subsubsection{Requisitos de Materiales}

\begin{itemize}
  \item Caja rectangular FG 4''x2'' con orejas para interruptor.
  \item Conectores tipo "romex" para ingreso de tuberia PVC.
  \item Tornillos de fijacion y tarugos de expansion para anclaje a muro.
  \item Cinta aislante vinilica.
\end{itemize}

\subsubsection{Metodo de Ejecucion}

\begin{enumerate}
  \item Replantear la ubicacion del interruptor en el muro segun el plano electrico, verificando que no quede detras de puertas o muebles fijos.
  \item Realizar la excavacion o corte en el muro para el empotramiento de la caja, con una profundidad suficiente para que la cara frontal quede al ras con la superficie del tarrajeo.
  \item Fijar la caja rectangular con tornillos y tarugos, verificando su nivelacion y plomo.
  \item Conectar la tuberia PVC de 3/4'' a la caja mediante conectores, asegurando la continuidad de la canalizacion.
  \item Dejar una longitud de conductor de 0.25 m desde el borde de la caja para conexiones futuras.
  \item Proteger la caja con tapa provisional durante la etapa de construccion.
\end{enumerate}

\subsubsection{Control de Calidad}

\begin{itemize}
  \item Verificar que la caja quede al ras con el muro terminado (tarrajeo).
  \item Comprobar la nivelacion vertical y horizontal de la caja.
  \item Inspeccionar que la tuberia ingrese correctamente a la caja mediante conectores.
  \item Confirmar la altura de instalacion contra el plano aprobado (1.20 m).
\end{itemize}

\subsubsection{Metodo de Medicion}

Se medira por pieza (pza) de caja rectangular para interruptor completamente instalada.

\subsubsection{Unidad de Medida}

pza (pieza).

\subsubsection{Forma de Pago}

Pago por pieza instalada y aprobada por la supervision, incluyendo materiales y mano de obra.

%------------------------------------------------------------------------------
\subsection{OE.5.2.1.2.1 Caja Rectangular FG 4''x2'' - Salida de Tomacorriente Doble con Tierra}

\subsubsection{Especificaciones Tecnicas}

La salida de tomacorriente doble con tierra consiste en la instalacion de una caja rectangular de fierro galvanizado (FG) de 4''x2'' para alojar un tomacorriente doble con conexion a tierra, de 15 A - 250 V. Se instalara empotrada en muros a una altura de 0.40 m sobre el nivel del piso terminado (en ambientes generales) y a 1.10 m en zonas de cocina sobre encimeras.

La instalacion debe incluir tres conductores: fase (rojo o marron), neutro (blanco o celeste) y proteccion a tierra (verde o verde-amarillo), todos de 2.5 mm2.

\subsubsection{Requisitos de Materiales}

\begin{itemize}
  \item Caja rectangular FG 4''x2'' con orejas para tomacorriente.
  \item Tomacorriente doble con tierra 15 A - 250 V, color blanco, marca TICINO o BTICINO.
  \item Placa frontal decorativa de 2 modulos, color blanco.
  \item Conector tipo "romex" para ingreso de tuberia PVC.
  \item Tornillos de fijacion y tarugos.
\end{itemize}

\subsubsection{Metodo de Ejecucion}

\begin{enumerate}
  \item Replantear la ubicacion del tomacorriente en el muro segun el plano electrico, respetando las alturas establecidas.
  \item Realizar el corte en el muro para empotrar la caja rectangular.
  \item Fijar la caja con tornillos y tarugos, asegurando nivelacion.
  \item Conectar la tuberia PVC de 3/4'' a la caja.
  \item Pasar los conductores de fase, neutro y tierra de 2.5 mm2 desde la caja de derivacion hasta la caja del tomacorriente.
  \item Instalar el tomacorriente doble conectando: fase al borne L, neutro al borne N y tierra al borne de tierra (símbolo $\perp$).
  \item Fijar la placa frontal decorativa.
\end{enumerate}

\subsubsection{Control de Calidad}

\begin{itemize}
  \item Verificar la polaridad correcta (fase en borne L, neutro en borne N).
  \item Comprobar la continuidad del conductor de proteccion a tierra.
  \item Medir la tension entre fase y neutro (debe ser 220 V ± 10\%).
  \item Inspeccionar la correcta fijacion de la placa frontal.
  \item Verificar que el tomacorriente funcione correctamente con un comprobador de tomacorrientes.
\end{itemize}

\subsubsection{Metodo de Medicion}

Se medira por pieza (pza) de tomacorriente doble con tierra completamente instalado, incluyendo caja, conductores, tomacorriente y placa frontal.

\subsubsection{Unidad de Medida}

pza (pieza).

\subsubsection{Forma de Pago}

Pago por pieza instalada, verificada con comprobador de tomacorrientes y aprobada por supervision.

%------------------------------------------------------------------------------
\subsection{OE.5.2.1.2.2 Caja Rectangular FG 4''x2'' - Salida de Tomacorriente Doble con Proteccion GFCI}

\subsubsection{Especificaciones Tecnicas}

La salida de tomacorriente con proteccion GFCI (Ground Fault Circuit Interrupter) consiste en la instalacion de un tomacorriente doble con interruptor diferencial incorporado, de 15 A - 250 V, con sensibilidad de 30 mA. Esta salida se instala en ambientes humedos como baños, cocinas y lavanderias, segun lo exigido por la Norma EM.010 y el CNE-U.

La caja rectangular FG 4''x2'' debe tener suficiente profundidad para alojar el mecanismo GFCI, que es de mayor volumen que un tomacorriente convencional.

\subsubsection{Requisitos de Materiales}

\begin{itemize}
  \item Caja rectangular FG 4''x2'' profunda (minimo 60 mm) para alojar mecanismo GFCI.
  \item Tomacorriente doble con proteccion GFCI 15 A - 250 V, 30 mA, marca TICINO o LEGRAND.
  \item Placa frontal decorativa de 2 modulos, color blanco.
  \item Conectores y accesorios de instalacion.
\end{itemize}

\subsubsection{Metodo de Ejecucion}

\begin{enumerate}
  \item Replantear ubicacion en el muro del ambiente humedo correspondiente.
  \item Instalar la caja rectangular empotrada en el muro siguiendo el mismo procedimiento del tomacorriente general.
  \\item Pasar los conductores de fase, neutro y tierra de 2.5 mm2.
  \item Conectar el tomacorriente GFCI: la fase al borne "LINE", el neutro al borne "LINE N", y la tierra al borne de tierra. Los bornes "LOAD" se utilizan solo si se requiere proteccion downstream.
  \item Instalar la placa frontal decorativa.
  \item Realizar prueba de funcionamiento presionando el boton "TEST" (debe disparar) y luego "RESET" (debe restablecer).
\end{enumerate}

\subsubsection{Control de Calidad}

\begin{itemize}
  \item Verificar la correcta conexion de los conductores segun el diagrama del fabricante.
  \item Realizar prueba de disparo con el boton TEST (debe cortar el suministro).
  \item Realizar prueba de restablecimiento con el boton RESET.
  \item Medir la corriente de fuga nominal si se dispone de instrumento adecuado.
  \item Verificar tension de salida 220 V ± 10\%.
\end{itemize}

\subsubsection{Metodo de Medicion}

Se medira por pieza (pza) de tomacorriente GFCI completamente instalado, incluyendo pruebas de funcionamiento.

\subsubsection{Unidad de Medida}

pza (pieza).

\subsubsection{Forma de Pago}

Pago por pieza instalada, probada (TEST/RESET verificado) y aprobada por supervision.

%------------------------------------------------------------------------------
\subsection{OE.5.2.1.2.3 Caja Rectangular FG 4''x2'' - Salida de Tomacorriente Servicio Pesado 20A}

\subsubsection{Especificaciones Tecnicas}

La salida de tomacorriente de servicio pesado consiste en la instalacion de un tomacorriente doble de 20 A - 250 V con conexion a tierra, disenado para cargas mayores como electrobombas, lavadoras o equipos de cocina de alta potencia. Se diferencia del tomacorriente estandar por su mayor capacidad de corriente y por admitir clavijas con espiga de tierra de mayor diametro.

Se instala en caja rectangular FG 4''x2'' con conductores de 2.5 mm2 minimo y proteccion termomagnetica de 20 A.

\subsubsection{Requisitos de Materiales}

\begin{itemize}
  \item Caja rectangular FG 4''x2''.
  \item Tomacorriente doble servicio pesado 20 A - 250 V con tierra, color blanco.
  \item Placa frontal decorativa de 2 modulos.
  \item Conductor de cobre LSOH 2.5 mm2 para fase, neutro y tierra.
  \item Accesorios de instalacion.
\end{itemize}

\subsubsection{Metodo de Ejecucion}

\begin{enumerate}
  \item Replantear ubicacion en el muro cerca del equipo a servir (electrobomba, lavadora).
  \item Instalar la caja rectangular empotrada en muro siguiendo el procedimiento estandar.
  \item Pasar conductores de 2.5 mm2 (fase, neutro y tierra).
  \item Conectar el tomacorriente de servicio pesado: fase al borne L, neutro al borne N, tierra al borne $\perp$.
  \item Fijar la placa frontal.
\end{enumerate}

\subsubsection{Control de Calidad}

\begin{itemize}
  \item Verificar que el tomacorriente soporte 20 A nominales (verificar sello o grabado).
  \item Comprobar polaridad y continuidad de tierra.
  \item Medir tension de salida 220 V ± 10\%.
  \item Verificar que la proteccion termomagnetica del circuito sea de 20 A.
\end{itemize}

\subsubsection{Metodo de Medicion}

Se medira por pieza (pza) de tomacorriente servicio pesado 20 A completamente instalado.

\subsubsection{Unidad de Medida}

pza (pieza).

\subsubsection{Forma de Pago}

Pago por pieza instalada, verificada y aprobada por supervision.
